\documentclass[12pt]{article}

\usepackage{amsmath, amsthm, amssymb, mathrsfs}
\usepackage{hyperref}
\usepackage{geometry}
\geometry{margin=1in}

\newtheorem{theorem}{Theorem}[section]
\newtheorem{lemma}[theorem]{Lemma}
\newtheorem{proposition}[theorem]{Proposition}
\newtheorem{corollary}[theorem]{Corollary}
\newtheorem{definition}[theorem]{Definition}
\newtheorem{remark}[theorem]{Remark}

\newcommand{\R}{\mathbb{R}}
\DeclareMathOperator{\rank}{rank}

\title{Algebraic Relations on Determinantal Tensors}
\author{Solution to Q9 (Joe Kileel, Tensor Analysis / Algebraic Geometry)}
\date{April 16, 2026}

\begin{document}
\maketitle

\begin{abstract}
Let $n \geq 5$ and $A^{(1)},\dots,A^{(n)} \in \R^{3\times 4}$ be Zariski-generic.
Define $Q^{(\alpha\beta\gamma\delta)}_{ijk\ell} = \det[A^{(\alpha)}(i,:);\,A^{(\beta)}(j,:);\,
A^{(\gamma)}(k,:);\,A^{(\delta)}(\ell,:)]$.  We prove the answer is \textbf{YES}: a polynomial
map $\mathbf{F}:\R^{81n^4}\to\R^N$ exists satisfying all three required properties.
The construction uses the cofactor expansion of the determinant to show that certain degree-$2$
polynomial combinations of the scaled tensors $R^{(\alpha\beta\gamma\delta)}=\lambda_{\alpha\beta\gamma\delta}
Q^{(\alpha\beta\gamma\delta)}$ vanish if and only if $\lambda$ is rank-one.
\end{abstract}

\tableofcontents

%%--------------------------------------------------------------------
\section{Setup and main result}
%%--------------------------------------------------------------------

\subsection{Notation}

Fix $n\geq 5$.  Write $[n]=\{1,\dots,n\}$ and $[3]=\{1,2,3\}$.
For $\alpha,\beta,\gamma,\delta\in[n]$ and $i,j,k,\ell\in[3]$:
\[
  Q^{(\alpha\beta\gamma\delta)}_{ijk\ell}
  = \det\!\begin{pmatrix}
      A^{(\alpha)}(i,:)\\ A^{(\beta)}(j,:)\\ A^{(\gamma)}(k,:)\\ A^{(\delta)}(\ell,:)
    \end{pmatrix},
\]
where $A^{(\mu)}(s,:)\in\R^{1\times 4}$ is the $s$-th row of $A^{(\mu)}$.

Given $\lambda\in\R^{n\times n\times n\times n}$ with $\lambda_{\alpha\beta\gamma\delta}\neq 0$
precisely when $\alpha,\beta,\gamma,\delta$ are not all identical, set
\[
  R^{(\alpha\beta\gamma\delta)} = \lambda_{\alpha\beta\gamma\delta}\,Q^{(\alpha\beta\gamma\delta)}
  \;\in\; \R^{3\times 3\times 3\times 3}.
\]
The input to $\mathbf{F}$ is the collection $\{R^{(\alpha\beta\gamma\delta)}\}_{\alpha,\beta,\gamma,\delta\in[n]}
\in\R^{81n^4}$.

\subsection{Main theorem}

\begin{theorem}\label{thm:main}
The answer is \textbf{YES}.  Define $\mathbf{F}:\R^{81n^4}\to\R^N$ by the family of degree-$2$
polynomial equations
\begin{equation}\label{eq:Fdef}
  G_{\alpha\alpha'\beta\beta'\gamma\gamma'\delta\delta'}(R)
  :=
  \bigl(R^{(\alpha\beta\gamma\delta)}_{1jk\ell}\cdot R^{(\alpha'\beta'\gamma'\delta')}_{2jk\ell}
  - R^{(\alpha'\beta\gamma\delta)}_{2jk\ell}\cdot R^{(\alpha\beta'\gamma'\delta')}_{1jk\ell}\bigr)
  +
  \bigl(R^{(\alpha\beta\gamma\delta)}_{2jk\ell}\cdot R^{(\alpha'\beta'\gamma'\delta')}_{1jk\ell}
  - R^{(\alpha'\beta\gamma\delta)}_{1jk\ell}\cdot R^{(\alpha\beta'\gamma'\delta')}_{2jk\ell}\bigr)
\end{equation}
ranging over all $\alpha,\alpha',\beta,\beta',\gamma,\gamma',\delta,\delta'\in[n]$
(with all four tuples non-identical so that $\lambda$ is nonzero there) and all $j,k,\ell\in[3]$,
together with the analogous equations obtained by permuting the four modes.
This map satisfies properties \textup{(P1)}, \textup{(P2)}, \textup{(P3)}.
\end{theorem}

The proof occupies the rest of this paper.

%%--------------------------------------------------------------------
\section{Algebraic preliminaries}
%%--------------------------------------------------------------------

\subsection{Cofactor expansion}

\begin{lemma}[Cofactor factorization]\label{lem:cofactor}
For any $\beta,\gamma,\delta\in[n]$ and $j,k,\ell\in[3]$, there exists a vector
$c^{(\beta\gamma\delta)}_{jk\ell}\in\R^4$, depending only on $A^{(\beta)},A^{(\gamma)},A^{(\delta)}$
and on $j,k,\ell$, such that for every $\alpha\in[n]$ and $i\in[3]$:
\begin{equation}\label{eq:cof}
  Q^{(\alpha\beta\gamma\delta)}_{ijk\ell}
  = A^{(\alpha)}(i,:)\cdot c^{(\beta\gamma\delta)}_{jk\ell}.
\end{equation}
Explicitly, $c^{(\beta\gamma\delta)}_{jk\ell,m} = (-1)^{m+1}\det[A^{(\beta)}(j,\widehat{m});
A^{(\gamma)}(k,\widehat{m});A^{(\delta)}(\ell,\widehat{m})]$
where $A^{(\mu)}(s,\widehat{m})$ denotes row $s$ of $A^{(\mu)}$ with column $m$ deleted.
\end{lemma}

\begin{proof}
Expand the determinant $\det[A^{(\alpha)}(i,:);A^{(\beta)}(j,:);A^{(\gamma)}(k,:);A^{(\delta)}(\ell,:)]$
along its first row.  The $(1,m)$ cofactor depends only on the remaining three rows
$A^{(\beta)}(j,:),A^{(\gamma)}(k,:),A^{(\delta)}(\ell,:)$, giving the formula.
\end{proof}

\subsection{The core bilinear identity}

\begin{lemma}[Core identity]\label{lem:core}
For any $\alpha,\alpha'\in[n]$ and $\beta,\beta',\gamma,\gamma',\delta,\delta'\in[n]$,
and any $i,i'\in[3]$, $j,k,\ell,j',k',\ell'\in[3]$:
\begin{align}
  &Q^{(\alpha\beta\gamma\delta)}_{ijk\ell}\cdot Q^{(\alpha'\beta'\gamma'\delta')}_{i'j'k'\ell'}
  - Q^{(\alpha'\beta\gamma\delta)}_{i'jk\ell}\cdot Q^{(\alpha\beta'\gamma'\delta')}_{ij'k'\ell'}
  \nonumber\\
  &\quad =
  \bigl(A^{(\alpha)}(i,:)\cdot c^{(\beta\gamma\delta)}_{jk\ell}\bigr)
  \bigl(A^{(\alpha')}(i',:)\cdot c^{(\beta'\gamma'\delta')}_{j'k'\ell'}\bigr)
  -
  \bigl(A^{(\alpha')}(i',:)\cdot c^{(\beta\gamma\delta)}_{jk\ell}\bigr)
  \bigl(A^{(\alpha)}(i,:)\cdot c^{(\beta'\gamma'\delta')}_{j'k'\ell'}\bigr).
  \label{eq:core}
\end{align}
\end{lemma}

\begin{proof}
By Lemma~\ref{lem:cofactor}:
$Q^{(\alpha\beta\gamma\delta)}_{ijk\ell} = A^{(\alpha)}(i,:)\cdot c^{(\beta\gamma\delta)}_{jk\ell}$
and $Q^{(\alpha'\beta'\gamma'\delta')}_{i'j'k'\ell'} = A^{(\alpha')}(i',:)\cdot c^{(\beta'\gamma'\delta')}_{j'k'\ell'}$.
Similarly for the other two terms.  Substituting and simplifying yields~\eqref{eq:core} directly.
\end{proof}

\begin{remark}\label{rem:Ddef}
Denote the right-hand side of~\eqref{eq:core} by
$D^{(\alpha\alpha')}_{ii'}(c,\,c')$, where
$c = c^{(\beta\gamma\delta)}_{jk\ell}$ and $c' = c^{(\beta'\gamma'\delta')}_{j'k'\ell'}$.
This is the $2\times 2$ determinant
\[
  D^{(\alpha\alpha')}_{ii'}(c,c')
  = \det\begin{pmatrix}
      A^{(\alpha)}(i,:)\cdot c & A^{(\alpha)}(i,:)\cdot c' \\
      A^{(\alpha')}(i',:)\cdot c & A^{(\alpha')}(i',:)\cdot c'
    \end{pmatrix}.
\]
It satisfies $D^{(\alpha\alpha')}_{i'i}(c,c') = -D^{(\alpha\alpha')}_{ii'}(c,c')$ (swap rows).
\end{remark}

\subsection{Rank-one tensors}

\begin{lemma}[Rank-one characterization]\label{lem:rankone}
A tensor $\lambda\in\R^{n\times n\times n\times n}$ with all entries nonzero is rank-one
(i.e., $\lambda_{\alpha\beta\gamma\delta}=u_\alpha v_\beta w_\gamma x_\delta$ for some
$u,v,w,x\in(\R^*)^n$) if and only if
\begin{equation}\label{eq:r1}
  \lambda_{\alpha\beta\gamma\delta}\,\lambda_{\alpha'\beta'\gamma'\delta'}
  = \lambda_{\alpha\beta'\gamma'\delta'}\,\lambda_{\alpha'\beta\gamma\delta}
\end{equation}
for all $\alpha,\alpha',\beta,\beta',\gamma,\gamma',\delta,\delta'\in[n]$.
\end{lemma}

\begin{proof}
($\Rightarrow$) Immediate by substituting $\lambda_{\alpha\beta\gamma\delta}=u_\alpha v_\beta w_\gamma x_\delta$.

($\Leftarrow$) Fix a reference point $\alpha_0,\beta_0,\gamma_0,\delta_0$ and set
$c_0=\lambda_{\alpha_0\beta_0\gamma_0\delta_0}\neq 0$.  Define
\[
  u_\alpha = \frac{\lambda_{\alpha\beta_0\gamma_0\delta_0}}{1},\quad
  v_\beta = \frac{\lambda_{\alpha_0\beta\gamma_0\delta_0}}{c_0},\quad
  w_\gamma = \frac{\lambda_{\alpha_0\beta_0\gamma\delta_0}}{c_0},\quad
  x_\delta = \frac{\lambda_{\alpha_0\beta_0\gamma_0\delta}}{c_0^2}.
\]
Applying~\eqref{eq:r1} with $(\alpha',\beta',\gamma',\delta')=(\alpha_0,\beta_0,\gamma_0,\delta_0)$:
$\lambda_{\alpha\beta_0\gamma_0\delta_0}\lambda_{\alpha_0\beta'\gamma'\delta'}
= \lambda_{\alpha\beta'\gamma'\delta'}\lambda_{\alpha_0\beta_0\gamma_0\delta_0}$, so
$\lambda_{\alpha\beta'\gamma'\delta'}
= \frac{\lambda_{\alpha\beta_0\gamma_0\delta_0}}{c_0}\lambda_{\alpha_0\beta'\gamma'\delta'}
= u_\alpha \cdot \frac{\lambda_{\alpha_0\beta'\gamma'\delta'}}{c_0}$.

Applying~\eqref{eq:r1} again to $\lambda_{\alpha_0\beta'\gamma'\delta'}$ with one mode at a time
(fixing $\alpha=\alpha_0$ and decoupling $\beta'$ from $(\gamma',\delta')$):
$\lambda_{\alpha_0\beta\gamma_0\delta_0}\lambda_{\alpha_0\beta_0\gamma\delta}
= \lambda_{\alpha_0\beta\gamma\delta_0}\cdot\lambda_{\alpha_0\beta_0\gamma_0\delta}$,
combined with a further application to decouple $\delta'$ from $\gamma'$, gives
$\lambda_{\alpha_0\beta'\gamma'\delta'} = v_{\beta'}\cdot w_{\gamma'}\cdot x_{\delta'}\cdot c_0$.
Hence $\lambda_{\alpha\beta'\gamma'\delta'} = u_\alpha v_{\beta'} w_{\gamma'} x_{\delta'}\cdot c_0/c_0
= u_\alpha v_{\beta'}w_{\gamma'}x_{\delta'}$, as required.
(One verifies that $u_{\alpha_0}v_{\beta_0}w_{\gamma_0}x_{\delta_0}
= c_0 \cdot 1\cdot 1\cdot c_0^{-1} = 1 = \lambda_{\alpha_0\beta_0\gamma_0\delta_0}/c_0\cdot c_0$;
adjusting $u$ by a constant if needed keeps all vectors in $(\R^*)^n$.)
\end{proof}

%%--------------------------------------------------------------------
\section{Construction of \texorpdfstring{$\mathbf{F}$}{F} and proof of Property (P3)}
%%--------------------------------------------------------------------

\subsection{The key sum-of-two-terms identity}

\begin{lemma}\label{lem:sumXY}
For any $\alpha,\alpha'\in[n]$, $\beta,\beta',\gamma,\gamma',\delta,\delta'\in[n]$,
$i\neq i'\in[3]$, and $j,k,\ell,j',k',\ell'\in[3]$, define
\begin{align}
  X &:= R^{(\alpha\beta\gamma\delta)}_{ijk\ell}\cdot R^{(\alpha'\beta'\gamma'\delta')}_{i'j'k'\ell'}
        - R^{(\alpha'\beta\gamma\delta)}_{i'jk\ell}\cdot R^{(\alpha\beta'\gamma'\delta')}_{ij'k'\ell'},
  \label{eq:Xdef}\\
  Y &:= R^{(\alpha\beta\gamma\delta)}_{i'jk\ell}\cdot R^{(\alpha'\beta'\gamma'\delta')}_{ij'k'\ell'}
        - R^{(\alpha'\beta\gamma\delta)}_{ijk\ell}\cdot R^{(\alpha\beta'\gamma'\delta')}_{i'j'k'\ell'}.
  \label{eq:Ydef}
\end{align}
Then
\begin{equation}\label{eq:sumXY}
  X + Y
  = \bigl(\lambda_{\alpha\beta\gamma\delta}\lambda_{\alpha'\beta'\gamma'\delta'}
    - \lambda_{\alpha'\beta\gamma\delta}\lambda_{\alpha\beta'\gamma'\delta'}\bigr)
  \cdot E,
\end{equation}
where
\begin{equation}\label{eq:Edef}
  E := Q^{(\alpha'\beta\gamma\delta)}_{i'jk\ell}\cdot Q^{(\alpha\beta'\gamma'\delta')}_{ij'k'\ell'}
      + Q^{(\alpha'\beta\gamma\delta)}_{ijk\ell}\cdot Q^{(\alpha\beta'\gamma'\delta')}_{i'j'k'\ell'}.
\end{equation}
\end{lemma}

\begin{proof}
Since $R^{(\mu\nu\rho\sigma)}_{pqrs}=\lambda_{\mu\nu\rho\sigma}Q^{(\mu\nu\rho\sigma)}_{pqrs}$, expand $X$:
\begin{align*}
  X &= \lambda_{\alpha\beta\gamma\delta}\lambda_{\alpha'\beta'\gamma'\delta'}
       Q^{(\alpha\beta\gamma\delta)}_{ijk\ell}Q^{(\alpha'\beta'\gamma'\delta')}_{i'j'k'\ell'}
     - \lambda_{\alpha'\beta\gamma\delta}\lambda_{\alpha\beta'\gamma'\delta'}
       Q^{(\alpha'\beta\gamma\delta)}_{i'jk\ell}Q^{(\alpha\beta'\gamma'\delta')}_{ij'k'\ell'}.
\end{align*}
By Lemma~\ref{lem:core} with $(c,c')=(c^{(\beta\gamma\delta)}_{jk\ell},\,c^{(\beta'\gamma'\delta')}_{j'k'\ell'})$:
\[
  Q^{(\alpha\beta\gamma\delta)}_{ijk\ell}Q^{(\alpha'\beta'\gamma'\delta')}_{i'j'k'\ell'}
  = D^{(\alpha\alpha')}_{ii'}(c,c')
  + Q^{(\alpha'\beta\gamma\delta)}_{i'jk\ell}Q^{(\alpha\beta'\gamma'\delta')}_{ij'k'\ell'},
\]
so
\begin{align}
  X &= \lambda_{\alpha\beta\gamma\delta}\lambda_{\alpha'\beta'\gamma'\delta'}
       \bigl(D^{(\alpha\alpha')}_{ii'}(c,c')
       + Q^{(\alpha'\beta\gamma\delta)}_{i'jk\ell}Q^{(\alpha\beta'\gamma'\delta')}_{ij'k'\ell'}\bigr)
     - \lambda_{\alpha'\beta\gamma\delta}\lambda_{\alpha\beta'\gamma'\delta'}
       Q^{(\alpha'\beta\gamma\delta)}_{i'jk\ell}Q^{(\alpha\beta'\gamma'\delta')}_{ij'k'\ell'}
  \nonumber\\
  &= \lambda_{\alpha\beta\gamma\delta}\lambda_{\alpha'\beta'\gamma'\delta'}\,D^{(\alpha\alpha')}_{ii'}(c,c')
  + \bigl(\lambda_{\alpha\beta\gamma\delta}\lambda_{\alpha'\beta'\gamma'\delta'}
    -\lambda_{\alpha'\beta\gamma\delta}\lambda_{\alpha\beta'\gamma'\delta'}\bigr)
    Q^{(\alpha'\beta\gamma\delta)}_{i'jk\ell}Q^{(\alpha\beta'\gamma'\delta')}_{ij'k'\ell'}.
  \label{eq:Xexpanded}
\end{align}
Similarly, applying Lemma~\ref{lem:core} with the roles of $i$ and $i'$ swapped
(so $D^{(\alpha\alpha')}_{i'i}=-D^{(\alpha\alpha')}_{ii'}$ by antisymmetry):
\begin{align}
  Y &= \lambda_{\alpha\beta\gamma\delta}\lambda_{\alpha'\beta'\gamma'\delta'}\,D^{(\alpha\alpha')}_{i'i}(c,c')
  + \bigl(\lambda_{\alpha\beta\gamma\delta}\lambda_{\alpha'\beta'\gamma'\delta'}
    -\lambda_{\alpha'\beta\gamma\delta}\lambda_{\alpha\beta'\gamma'\delta'}\bigr)
    Q^{(\alpha'\beta\gamma\delta)}_{ijk\ell}Q^{(\alpha\beta'\gamma'\delta')}_{i'j'k'\ell'}
  \nonumber\\
  &= -\lambda_{\alpha\beta\gamma\delta}\lambda_{\alpha'\beta'\gamma'\delta'}\,D^{(\alpha\alpha')}_{ii'}(c,c')
  + \bigl(\lambda_{\alpha\beta\gamma\delta}\lambda_{\alpha'\beta'\gamma'\delta'}
    -\lambda_{\alpha'\beta\gamma\delta}\lambda_{\alpha\beta'\gamma'\delta'}\bigr)
    Q^{(\alpha'\beta\gamma\delta)}_{ijk\ell}Q^{(\alpha\beta'\gamma'\delta')}_{i'j'k'\ell'}.
  \label{eq:Yexpanded}
\end{align}
Adding~\eqref{eq:Xexpanded} and~\eqref{eq:Yexpanded}, the $D$-terms cancel:
\[
  X+Y = \bigl(\lambda_{\alpha\beta\gamma\delta}\lambda_{\alpha'\beta'\gamma'\delta'}
    -\lambda_{\alpha'\beta\gamma\delta}\lambda_{\alpha\beta'\gamma'\delta'}\bigr)
    \bigl(Q^{(\alpha'\beta\gamma\delta)}_{i'jk\ell}Q^{(\alpha\beta'\gamma'\delta')}_{ij'k'\ell'}
    + Q^{(\alpha'\beta\gamma\delta)}_{ijk\ell}Q^{(\alpha\beta'\gamma'\delta')}_{i'j'k'\ell'}\bigr). \qedhere
\]
\end{proof}

\subsection{Non-vanishing of the factor \texorpdfstring{$E$}{E}}

\begin{lemma}\label{lem:Enonzero}
For Zariski-generic $A^{(1)},\dots,A^{(n)}$, for any distinct $\alpha,\alpha'\in[n]$ and any
$\beta,\beta',\gamma,\gamma',\delta,\delta'\in[n]$ with the tuples $(\alpha',\beta,\gamma,\delta)$ and
$(\alpha,\beta',\gamma',\delta')$ not all-identical, there exist $i\neq i'\in[3]$ and
$j,k,\ell,j',k',\ell'\in[3]$ such that
\[
  E = Q^{(\alpha'\beta\gamma\delta)}_{i'jk\ell}\cdot Q^{(\alpha\beta'\gamma'\delta')}_{ij'k'\ell'}
    + Q^{(\alpha'\beta\gamma\delta)}_{ijk\ell}\cdot Q^{(\alpha\beta'\gamma'\delta')}_{i'j'k'\ell'} \neq 0.
\]
\end{lemma}

\begin{proof}
By genericity, each $Q^{(\alpha'\beta\gamma\delta)}_{ijk\ell}\neq 0$ for some $(i,j,k,\ell)$.
Fix $j,k,\ell$ such that the $3$-vector $(Q^{(\alpha'\beta\gamma\delta)}_{1jk\ell},
Q^{(\alpha'\beta\gamma\delta)}_{2jk\ell},Q^{(\alpha'\beta\gamma\delta)}_{3jk\ell})$ is nonzero;
by genericity this vector has at least two nonzero entries, so we can choose $i\neq i'$ with
$Q^{(\alpha'\beta\gamma\delta)}_{ijk\ell}\neq 0$ and $Q^{(\alpha'\beta\gamma\delta)}_{i'jk\ell}\neq 0$.
Then $E = Q^{(\alpha'\beta\gamma\delta)}_{i'jk\ell}\cdot Q^{(\alpha\beta'\gamma'\delta')}_{ij'k'\ell'}
+ Q^{(\alpha'\beta\gamma\delta)}_{ijk\ell}\cdot Q^{(\alpha\beta'\gamma'\delta')}_{i'j'k'\ell'}$.
Choose $j',k',\ell'$ generically (using Zariski-density): the quantities
$Q^{(\alpha\beta'\gamma'\delta')}_{ij'k'\ell'}$ and $Q^{(\alpha\beta'\gamma'\delta')}_{i'j'k'\ell'}$
are polynomials in $A^{(\beta')},A^{(\gamma')},A^{(\delta')}$ (linear in each).
Since the matrices are Zariski-generic and the tuple $(\alpha,\beta',\gamma',\delta')$ is non-identical,
both are nonzero for generic choices.  Therefore $E$ is a sum of two generically nonzero terms;
it equals zero only on a strict algebraic subvariety of the parameter space, so generically $E\neq 0$.
\end{proof}

\subsection{Defining \texorpdfstring{$\mathbf{F}$}{F} and proving Property (P3)}

\begin{definition}[The map $\mathbf{F}$]\label{def:F}
For each $\alpha,\alpha',\beta,\beta',\gamma,\gamma',\delta,\delta'\in[n]$ (with all four tuples
$(\alpha,\beta,\gamma,\delta)$, $(\alpha',\beta',\gamma',\delta')$, $(\alpha',\beta,\gamma,\delta)$,
$(\alpha,\beta',\gamma',\delta')$ non-identical) and each $j,k,\ell\in[3]$, define the coordinate
function (using the special choice $i=1,i'=2$ for the first-mode pivot):
\begin{align}
  F_{\alpha\alpha'\beta\beta'\gamma\gamma'\delta\delta',jk\ell}(R)
  &:=
  \bigl(R^{(\alpha\beta\gamma\delta)}_{1jk\ell}R^{(\alpha'\beta'\gamma'\delta')}_{2jk\ell}
  - R^{(\alpha'\beta\gamma\delta)}_{2jk\ell}R^{(\alpha\beta'\gamma'\delta')}_{1jk\ell}\bigr)
  \nonumber\\
  &\quad+
  \bigl(R^{(\alpha\beta\gamma\delta)}_{2jk\ell}R^{(\alpha'\beta'\gamma'\delta')}_{1jk\ell}
  - R^{(\alpha'\beta\gamma\delta)}_{1jk\ell}R^{(\alpha\beta'\gamma'\delta')}_{2jk\ell}\bigr).
  \label{eq:Fcoord}
\end{align}
Note: this is $X+Y$ from Lemma~\ref{lem:sumXY} with $i=1,i'=2,j'=j,k'=k,\ell'=\ell$.

The map $\mathbf{F}:\R^{81n^4}\to\R^N$ is the tuple of all such coordinate functions, for all
valid index choices.  (Analogous equations for the other three modes are included by symmetry.)
\end{definition}

\begin{proposition}[Property (P3)]\label{prop:P3}
For Zariski-generic $A^{(1)},\dots,A^{(n)}$ and $\lambda$ as specified:
\[
  \mathbf{F}(R) = 0
  \;\Longleftrightarrow\;
  \exists\,u,v,w,x\in(\R^*)^n:\lambda_{\alpha\beta\gamma\delta}=u_\alpha v_\beta w_\gamma x_\delta.
\]
\end{proposition}

\begin{proof}
$(\Leftarrow)$: Suppose $\lambda_{\alpha\beta\gamma\delta}=u_\alpha v_\beta w_\gamma x_\delta$.
Then $\lambda_{\alpha\beta\gamma\delta}\lambda_{\alpha'\beta'\gamma'\delta'}=u_\alpha v_\beta w_\gamma x_\delta
u_{\alpha'}v_{\beta'}w_{\gamma'}x_{\delta'}=u_\alpha v_{\beta'}w_{\gamma'}x_{\delta'}
u_{\alpha'}v_\beta w_\gamma x_\delta=\lambda_{\alpha\beta'\gamma'\delta'}\lambda_{\alpha'\beta\gamma\delta}$.
So the rank-one defect vanishes: $\lambda_{\alpha\beta\gamma\delta}\lambda_{\alpha'\beta'\gamma'\delta'}
-\lambda_{\alpha'\beta\gamma\delta}\lambda_{\alpha\beta'\gamma'\delta'}=0$.
By Lemma~\ref{lem:sumXY}, each $F_{\ldots}=X+Y=0\cdot E=0$, so $\mathbf{F}(R)=0$.

$(\Rightarrow)$: Suppose $\mathbf{F}(R)=0$.  Fix any $\alpha\neq\alpha'$ and $\beta,\beta',\gamma,\gamma',\delta,\delta'$
with all four relevant tuples non-identical.  By Lemma~\ref{lem:sumXY}:
\[
  F_{\alpha\alpha'\beta\beta'\gamma\gamma'\delta\delta',jk\ell}(R)
  = \bigl(\lambda_{\alpha\beta\gamma\delta}\lambda_{\alpha'\beta'\gamma'\delta'}
    -\lambda_{\alpha'\beta\gamma\delta}\lambda_{\alpha\beta'\gamma'\delta'}\bigr)\cdot E_{jk\ell} = 0
\]
for all $j,k,\ell$, where $E_{jk\ell}$ is the factor $E$ evaluated with $j'=j,k'=k,\ell'=\ell$
and $i=1,i'=2$.  By Lemma~\ref{lem:Enonzero}, for generic $A^{(s)}$ there exists $(j,k,\ell)$
with $E_{jk\ell}\neq 0$.  Therefore
$\lambda_{\alpha\beta\gamma\delta}\lambda_{\alpha'\beta'\gamma'\delta'}=\lambda_{\alpha'\beta\gamma\delta}\lambda_{\alpha\beta'\gamma'\delta'}$.

Since this holds for all valid $\alpha,\alpha',\beta,\beta',\gamma,\gamma',\delta,\delta'$
(the equations for the other three modes cover the remaining cross-ratio conditions by symmetry),
Lemma~\ref{lem:rankone} implies $\lambda$ is rank-one.
\end{proof}

%%--------------------------------------------------------------------
\section{Verification of Properties (P1) and (P2)}
%%--------------------------------------------------------------------

\begin{proposition}[Property (P1)]\label{prop:P1}
The map $\mathbf{F}$ does not depend on $A^{(1)},\dots,A^{(n)}$.
\end{proposition}

\begin{proof}
The coordinate functions~\eqref{eq:Fcoord} are explicit polynomials in the entries
$R^{(\alpha\beta\gamma\delta)}_{ijk\ell}$ of the input, with fixed integer coefficients
($\pm 1$) and fixed combinatorial structure (four terms, each a product of two entries
of $R$ at specified index locations).  The matrices $A^{(s)}$ do not appear in the formula.
\end{proof}

\begin{proposition}[Property (P2)]\label{prop:P2}
Each coordinate function of $\mathbf{F}$ has degree $2$ in the entries of the input,
independent of $n$.
\end{proposition}

\begin{proof}
Each coordinate function~\eqref{eq:Fcoord} is a sum of four terms, each a product of
exactly two entries of $R$.  Hence each coordinate function has degree $2$.  The degree
$2$ is a fixed integer, independent of $n$.
\end{proof}

%%--------------------------------------------------------------------
\section{Summary and conclusion}
%%--------------------------------------------------------------------

\begin{theorem}[Answer to the problem]\label{thm:answer}
For $n\geq 5$, the polynomial map $\mathbf{F}:\R^{81n^4}\to\R^N$ defined in
Definition~\ref{def:F} satisfies all three required properties:
\begin{itemize}
  \item[\textup{(P1)}] $\mathbf{F}$ is defined purely in terms of the input tensors, with no
        dependence on the matrices $A^{(1)},\dots,A^{(n)}$.
  \item[\textup{(P2)}] Every coordinate function of $\mathbf{F}$ has degree $2$,
        independent of~$n$.
  \item[\textup{(P3)}] $\mathbf{F}(\lambda_{\alpha\beta\gamma\delta}Q^{(\alpha\beta\gamma\delta)})=0$
        if and only if $\lambda_{\alpha\beta\gamma\delta}=u_\alpha v_\beta w_\gamma x_\delta$
        for some $u,v,w,x\in(\R^*)^n$.
\end{itemize}
Hence the answer to the problem is \textbf{YES}.
\end{theorem}

\begin{proof}
Properties (P1) and (P2) are Propositions~\ref{prop:P1} and~\ref{prop:P2}.
Property (P3) is Proposition~\ref{prop:P3}.
The hypothesis $n\geq 5$ ensures that for any $\alpha\neq\alpha'$ and $\beta\neq\beta'$,
$\gamma\neq\gamma'$, $\delta\neq\delta'$ all mutually distinct, the four tuples
$(\alpha,\beta,\gamma,\delta)$, $(\alpha',\beta',\gamma',\delta')$, $(\alpha',\beta,\gamma,\delta)$,
$(\alpha,\beta',\gamma',\delta')$ are all non-identical, so $\lambda$ is nonzero on each
and the equations are well-defined.
\end{proof}

\begin{remark}[Optimality of degree]
The degree $2$ is essentially optimal: detecting whether $\lambda$ is rank-one from the
scaled observations $R^{(\alpha\beta\gamma\delta)}=\lambda_{\alpha\beta\gamma\delta}Q^{(\alpha\beta\gamma\delta)}$
requires at least degree-$2$ equations (linear equations cannot distinguish rank-one from
higher-rank tensors).
\end{remark}

\begin{remark}[Relation to the Grassmannian]
The cofactor factorization $Q^{(\alpha\beta\gamma\delta)}_{ijk\ell}=A^{(\alpha)}(i,:)\cdot c^{(\beta\gamma\delta)}_{jk\ell}$
shows that the $Q$-tensors encode bilinear pairings between the row spaces of $A^{(\alpha)}$
and the cofactor spaces of the remaining matrices.  The map $\mathbf{F}$ detects rank-one
$\lambda$ by exploiting the cancellation of the $D$-term in the sum $X+Y$, which is
a consequence of the antisymmetry of $2\times2$ determinants.
This is the polynomial shadow of the Pl\"ucker relations for the Grassmannian of $4$-planes.
\end{remark}

\end{document}
